\documentclass[10pt,a4paper]{amsart}
\usepackage{amsmath,amssymb,amsthm,graphicx}
% Retained minimal preamble for the Stacks Project chapter more-algebra.tex (More on Algebra).
% The chapter file pinned-original-source.tex carries only \input{preamble} on its first line, so no
% preamble text lives inside the chapter itself.  The Stacks Project uses one single shared preamble for
% every chapter; the snapshot stages it identically as sources/stacks-chow/preamble.tex and
% sources/stacks-etale-cohomology/preamble.tex.  The block below is copied mechanically, line for line,
% from sources/stacks-chow/preamble.tex lines 14-18 (the Xy-pic setup) and 154-217 (the native theorem
% declarations together with the mathematical macros).  Nothing was added, removed or redefined; the
% wrapper class and the amsmath/amssymb/amsthm/graphicx base packages are supplied by the assembler.
% Chapter: more-algebra.tex -- More on Algebra.
% For commutative diagrams we use Xy-pic
\usepackage[all]{xy}

% We use 2cell for 2-commutative diagrams.
\xyoption{2cell}

% Theorem environments.
%
\theoremstyle{plain}
\newtheorem{theorem}[subsection]{Theorem}
\newtheorem{proposition}[subsection]{Proposition}
\newtheorem{lemma}[subsection]{Lemma}

\theoremstyle{definition}
\newtheorem{definition}[subsection]{Definition}
\newtheorem{example}[subsection]{Example}
\newtheorem{exercise}[subsection]{Exercise}
\newtheorem{situation}[subsection]{Situation}

\theoremstyle{remark}
\newtheorem{remark}[subsection]{Remark}
\newtheorem{remarks}[subsection]{Remarks}

\numberwithin{equation}{subsection}

% Macros
%
\def\lim{\mathop{\mathrm{lim}}\nolimits}
\def\colim{\mathop{\mathrm{colim}}\nolimits}
\def\Spec{\mathop{\mathrm{Spec}}}
\def\Hom{\mathop{\mathrm{Hom}}\nolimits}
\def\Ext{\mathop{\mathrm{Ext}}\nolimits}
\def\SheafHom{\mathop{\mathcal{H}\!\mathit{om}}\nolimits}
\def\SheafExt{\mathop{\mathcal{E}\!\mathit{xt}}\nolimits}
\def\Sch{\mathit{Sch}}
\def\Mor{\mathop{\mathrm{Mor}}\nolimits}
\def\Ob{\mathop{\mathrm{Ob}}\nolimits}
\def\Sh{\mathop{\mathit{Sh}}\nolimits}
\def\NL{\mathop{N\!L}\nolimits}
\def\CH{\mathop{\mathrm{CH}}\nolimits}
\def\proetale{{pro\text{-}\acute{e}tale}}
\def\etale{{\acute{e}tale}}
\def\QCoh{\mathit{QCoh}}
\def\Ker{\mathop{\mathrm{Ker}}}
\def\Im{\mathop{\mathrm{Im}}}
\def\Coker{\mathop{\mathrm{Coker}}}
\def\Coim{\mathop{\mathrm{Coim}}}

% Boxtimes
%
\DeclareMathSymbol{\boxtimes}{\mathbin}{AMSa}{"02}

%
% Macros for moduli stacks/spaces
%
\def\QCohstack{\mathcal{QC}\!\mathit{oh}}
\def\Cohstack{\mathcal{C}\!\mathit{oh}}
\def\Spacesstack{\mathcal{S}\!\mathit{paces}}
\def\Quotfunctor{\mathrm{Quot}}
\def\Hilbfunctor{\mathrm{Hilb}}
\def\Curvesstack{\mathcal{C}\!\mathit{urves}}
\def\Polarizedstack{\mathcal{P}\!\mathit{olarized}}
\def\Complexesstack{\mathcal{C}\!\mathit{omplexes}}
% \Pic is the operator that assigns to X its picard group, usage \Pic(X)
% \Picardstack_{X/B} denotes the Picard stack of X over B
% \Picardfunctor_{X/B} denotes the Picard functor of X over B
\def\Pic{\mathop{\mathrm{Pic}}\nolimits}
\def\Picardstack{\mathcal{P}\!\mathit{ic}}
\def\Picardfunctor{\mathrm{Pic}}
\def\Deformationcategory{\mathcal{D}\!\mathit{ef}}

% --- appended 2026-10-12 (Chief): the project own definitions of the reference/slogan/history
% display environments, copied mechanically from the shared Stacks preamble lines 8-12
% (\usepackage{verbatim} plus the three \newenvironment lines, which the chapter preamble keeps
% verbatim from the project preamble). Nothing added, removed or redefined: these environments
% carry no mathematical content and the project itself typesets them as comments. Without these
% three lines a counted statement whose own text uses them fails with Environment slogan undefined.
\usepackage{verbatim}
\newenvironment{reference}{\comment}{\endcomment}
\newenvironment{slogan}{\comment}{\endcomment}
\newenvironment{history}{\comment}{\endcomment}
\usepackage{cleveref}
\crefname{originallemma}{original lemma}{original lemmas}
\Crefname{originallemma}{Original lemma}{Original lemmas}
\makeatletter
\newlabel{more-algebra-lemma-check-isomorphism-via-depth-and-ass}{{0AV8}{source}{}{sourcepointer-1}{}}
\newlabel{more-algebra-lemma-check-isomorphism-via-depth-and-ass@cref}{{[originallemma][0][]0AV8}{[1][0][]source}{}{}{}}
\makeatother
\crefname{originallemma}{original lemma}{original lemmas}
\Crefname{originallemma}{Original lemma}{Original lemmas}
\makeatletter
\newlabel{lemma-check-injective-on-ass}{{0AV7}{source}{}{sourcepointer-2}{}}
\newlabel{lemma-check-injective-on-ass@cref}{{[originallemma][0][]0AV7}{[1][0][]source}{}{}{}}
\makeatother
\crefname{originallemma}{original lemma}{original lemmas}
\Crefname{originallemma}{Original lemma}{Original lemmas}
\makeatletter
\newlabel{algebra-lemma-ideal-nonzerodivisor}{{00LL}{source}{}{sourcepointer-3}{}}
\newlabel{algebra-lemma-ideal-nonzerodivisor@cref}{{[originallemma][0][]00LL}{[1][0][]source}{}{}{}}
\makeatother
\crefname{originallemma}{original lemma}{original lemmas}
\Crefname{originallemma}{Original lemma}{Original lemmas}
\makeatletter
\newlabel{algebra-lemma-depth-in-ses}{{00LX}{source}{}{sourcepointer-4}{}}
\newlabel{algebra-lemma-depth-in-ses@cref}{{[originallemma][0][]00LX}{[1][0][]source}{}{}{}}
\makeatother
\crefname{originallemma}{original lemma}{original lemmas}
\Crefname{originallemma}{Original lemma}{Original lemmas}
\makeatletter
\newlabel{algebra-lemma-ass-zero}{{0587}{source}{}{sourcepointer-5}{}}
\newlabel{algebra-lemma-ass-zero@cref}{{[originallemma][0][]0587}{[1][0][]source}{}{}{}}
\makeatother
\title{A criterion for a map of modules to be an isomorphism in terms of depth and associated primes}
\author{The Stacks Project Authors}
\date{stacks-project:more-algebra.tex:local-snapshot-20261012; GNU FDL 1.2 or later}
\begin{document}
\maketitle
\setcounter{subsection}{1}
\noindent\textit{Editorial note.} All delivered mathematical text is the author's own, copied from the pinned Stacks Project chapter file more-algebra.tex; no statement or proof step was rewritten, reordered or removed.\par\medskip
\noindent\textit{Editorial note.} The chapter file pinned-original-source.tex carries only \textbackslash{}input\{preamble\} on its first line, so no preamble text lives inside the chapter. The retained file sources/stacks-more-algebra/retained-preamble.tex is a mechanical, line-for-line copy of the project's single shared preamble as staged in the snapshot (Xy-pic setup plus theorem declarations and mathematical macros). Nothing was added, removed or redefined, and the excerpt serves only as preamble, never as a span coordinate.\par\medskip
\noindent\textit{Editorial note.} Both counted spans start at the \textbackslash{}begin line of their environment, so the statement and the author's proof are delivered whole; surrounding prose of the section is deliberately outside the counted spans and is carried as standard prerequisites instead.\par\medskip
\noindent\textit{Editorial note.} The citation fields (source\_id, source\_version, authors, source\_url, license.name, source\_title) are identical to those of every other card this lane has produced from this chapter, so the batch selector treats them as one source.\par\medskip
\noindent\textit{Allowed standard prerequisites (editorial).}\begin{itemize}
\item Injectivity criterion on associated primes. Lemma \textbackslash{}ref\{lemma-check-injective-on-ass\}, cited at proof lines 5466-5467. Gives injectivity of the map under the stated hypotheses.
\item The associated primes of a localisation and the depth of a localisation. Algebra, Lemma \textbackslash{}ref\{algebra-lemma-ideal-nonzerodivisor\}, cited at proof lines 5477-5479. Turns the condition that a prime is not associated into a lower bound on the depth of the target.
\item Behaviour of depth in a short exact sequence. Algebra, Lemma \textbackslash{}ref\{algebra-lemma-depth-in-ses\}, cited at proof lines 5483-5485. Shows the cokernel has depth at least one at every prime where it need not already vanish.
\item A module with no associated primes is zero. Algebra, Lemma \textbackslash{}ref\{algebra-lemma-ass-zero\}, cited at proof lines 5485-5486. Concludes that the cokernel vanishes.
\end{itemize}\medskip
\ifdefined\hypertarget\hypertarget{sourcepointer-1}{}\fi
\noindent\textit{Source reference.} Original lemma 0AV8: focused-seeds/AGP/sources/stacks-snapshot-20261011/tags.txt, source lines 13676--13676. The lemma's own canonical tag, quoted so that readers can resolve the label inside the delivered text; it points at the counted statement itself and carries no proof content.\par\medskip
\ifdefined\hypertarget\hypertarget{sourcepointer-2}{}\fi
\noindent\textit{Source reference.} Original lemma 0AV7: focused-seeds/AGP/sources/stacks-snapshot-20261011/tags.txt, source lines 13675--13675. Navigation pointer to a statement the counted proof refers to but does not contain; nothing load-bearing is delegated to it.\par\medskip
\ifdefined\hypertarget\hypertarget{sourcepointer-3}{}\fi
\noindent\textit{Source reference.} Original lemma 00LL: focused-seeds/AGP/sources/stacks-snapshot-20261011/tags.txt, source lines 791--791. Navigation pointer to a statement the counted proof refers to but does not contain; nothing load-bearing is delegated to it.\par\medskip
\ifdefined\hypertarget\hypertarget{sourcepointer-4}{}\fi
\noindent\textit{Source reference.} Original lemma 00LX: focused-seeds/AGP/sources/stacks-snapshot-20261011/tags.txt, source lines 803--803. Navigation pointer to a statement the counted proof refers to but does not contain; nothing load-bearing is delegated to it.\par\medskip
\ifdefined\hypertarget\hypertarget{sourcepointer-5}{}\fi
\noindent\textit{Source reference.} Original lemma 0587: focused-seeds/AGP/sources/stacks-snapshot-20261011/tags.txt, source lines 6770--6770. Navigation pointer to a statement the counted proof refers to but does not contain; nothing load-bearing is delegated to it.\par\medskip
\begin{lemma}
\label{lemma-check-isomorphism-via-depth-and-ass}
Let $R$ be a Noetherian ring. Let $\varphi : M \to N$ be a map of
$R$-modules. Assume $M$ is finite and that for every prime $\mathfrak p$
of $R$ one of the following happens
\begin{enumerate}
\item $M_\mathfrak p \to N_\mathfrak p$ is an isomorphism, or
\item $\text{depth}(M_\mathfrak p) \geq 2$ and
$\mathfrak p \not \in \text{Ass}(N)$.
\end{enumerate}
Then $\varphi$ is an isomorphism.
\end{lemma}

\begin{proof}
By Lemma \ref{lemma-check-injective-on-ass} we see that $\varphi$ is injective.
Let $N' \subset N$ be an finitely generated $R$-module containing
the image of $M$. Then $\text{Ass}(N_\mathfrak p) = \emptyset$ implies
$\text{Ass}(N'_\mathfrak p) = \emptyset$.
Hence the assumptions of the lemma hold for $M \to N'$.
In order to prove that $\varphi$ is an isomorphism, it suffices
to prove the same thing for every such $N' \subset N$. Thus we may
assume $N$ is a finite $R$-module. In this case,
$\mathfrak p \not \in \text{Ass}(N) \Rightarrow
\text{depth}(N_\mathfrak p) \geq 1$, see
Algebra, Lemma \ref{algebra-lemma-ideal-nonzerodivisor}.
Consider the short exact sequence
$$
0 \to M \to N \to Q \to 0
$$
defining $Q$.
Looking at the conditions we see that either $Q_\mathfrak p = 0$
in case (1) or $\text{depth}(Q_\mathfrak p) \geq 1$ in case (2)
by Algebra, Lemma \ref{algebra-lemma-depth-in-ses}.
This implies that $Q$ does not have any associated primes, hence $Q = 0$ by
Algebra, Lemma \ref{algebra-lemma-ass-zero}.
\end{proof}

\end{document}
